
\section{Case studies of democratization}\label{section:case_studies}

\subsection{Sweden, 1917--1924}\label{section:SWEDEN}

The fall of the monarchy in Sweden offers an excellent example of a democratization associated with a large stock market response combined with subsequent redistribution. Relative to its Scandinavian neighbors, Sweden was slow to democratize. This changed in 1917. The year began with a conservative government in power. By autumn, however, this government had been forced from office due to ``food riots and the unreliability of the army'' \citep{Luebbert1991}. Worker and soldier unrest continued into 1918 and by October the decisive democratic breakthrough had occurred.  This victory brought with it a coalition Liberal-Social Democrat government from 1918--1920 which instituted several pro-labor policies through strengthening the already strong trade unions and instituting the 8 hour work day \citep{Bengtsson2014}. Universal suffrage was also established during this time, with the first elections under universal suffrage taking place between September 10th and 26th in 1921. V-Dem's Electoral Democracy Index tracks this progress well, as shown in Figure~\ref{fig:SWEDEN}, showing an initial increase in 1918 and final increase in 1922 as the newly elected government takes power.

While these policy changes did not immediately bring forth the famed Swedish welfare state---that would come about during and after the Great Depression---they did alter the bargaining power between labor and capital tremendously. For example, \cite{Bengtsson2014} finds a structural break in the capital share of income in 1920, with the capital share going from a high of 40\% in 1916 down to 20\% just after 1920. Moreover, this effect seemed to permanent; from 1920--2000, it would not reach above 30\% again.

\begin{figure}[t!]
    \captionsetup{width=1\linewidth}
	\caption{\textbf{Electoral Democracy Index and dividend yield, Sweden 1917--1924}}
	\renewcommand{\baselinestretch}{1}
    \captionsetup{singlelinecheck=off}
	\label{fig:SWEDEN}
	\footnotesize 
	\vspace{-.75cm}
    \begin{center}
		\includegraphics[width=1\textwidth]{../figures/raw/figure_F12_sweden_case_study.pdf}
	\end{center}
\end{figure}

Additional support for a nearly immediate reduction in inequality comes from examining top income shares. While exact numbers on how much inequality declined after the democratization are somewhat contested, recent research by \cite*{Bengtsson2021} on a random sample of tax returns in Stockholm indicate that the Gini coefficient fell by at much as 20 percentage points and the top 10\% share of income by 15 percentage points from 1920 to 1940. For comparison, the World Inequality Database (WID) reports that the top 10\% share in the United Kingdom and United States remained flat over this period, and in France only declined by 5 percentage points.  Similarly, the WID reports that the top 1\% income share in Sweden fell by 8 percentage points, compared to 5 percentage points in the UK and France from 1919--1941. It remained flat in the U.S. over this period. \cite{Bengtsson2019} also notes the discontinuity in Swedish income inequality post-democratization.

Finally, the Swedish democratization brought with it large increases in the dividend yield as shown in Figure~\ref{fig:SWEDEN}. The dividend yield began to rise in 1917 with the labor unrest and calls for increased political rights. In the year of the democratic breakthrough, the dividend yield rose further with the onset of the democratic breakthrough. From 1917--1920, the outcome of the democratization remained highly uncertain. However, as 1920 came to an end, the shift of power toward the left became complete, and brought with it large declines to inequality. With this uncertainty resolved, the dividend yield began to fall to its pre-democratization levels.

\subsection{France, 1847--1848}\label{section:FRANCE}

The establishing of the Second French Republic in the wake of the revolution of 1848 presents an excellent example of a failed democratization.  The movement toward 1848 began in 1847 with the beginning of the Reformist ``banquets'' at which toasts were drunk to the \textit{R\'{e}publique fran\c{c}aise} \citep{Marx1850}. This \textit{Campagne des banquets} was constructed to circumvent the restriction on political gatherings levied by the monarchy. While mostly liberal in nature, these banquets were also attended by reformists of all kinds; for example, a young Friedrich Engels attended some of these banquets starting in October 1847. King Louis Philippe allowed for these Reformist meetings to continue, resulting in an increase in free expression in the Electoral Democracy Index, as shown in Figure~\ref{fig:FRANCE}.  

As the banquets became more revolutionary in nature,  however, the Prime Minister of France, Fran\c{c}ois Guizot, outlawed them in January, 1848. Despite this ban, the gatherings continued. Things came to a head on February 22nd, when the French government banned the banquets for the second time, leading the organizing committee to cancel the events. Workers and students, however, had been mobilizing prior to the ban, and they did not plan to cancel their demonstration. It was with these demonstrations that a second ``Three Glorious Days'' began, leading to the ousting of King Louis Philippe on February 24th.

Shortly after the abdication of King Louis Philippe, the Second French Republic was declared. However, the democratic progress was short lived. Infighting in the proto-socialist groups made them politically ineffectual and ultimately led to the election of Louis Napoleon Bonaparte in the election of 1848. Bonaparte, a man viewed as the arch-ally to the \textit{bourgeoisie} by Marx, ultimately fully reversed the democratic progress in his famed 1851 \textit{coup d'\'{e}tat}, which established the Second French Empire.

Also shown in Table~\ref{fig:FRANCE} is the movement in the dividend yield across the failed democratization. Dividend yields spike in 1848 with the initial unrest and fall of the monarch. They then drop after the election of Louis Napoleon, but remain elevated until 1851, and the establishment of the Second Empire. In 1851 and 1852, stock prices rose 27\% and 53\%, respectively, signaling both the end of the episode, and investors' satisfaction with the 1851 coup.

\begin{figure}[t!]
    \captionsetup{width=1\linewidth}
	\caption{\textbf{Electoral Democracy Index and dividend yield, France 1847--1848}}
	\renewcommand{\baselinestretch}{1}
    \captionsetup{singlelinecheck=off}
	\label{fig:FRANCE}
	\footnotesize 
	\vspace{-.75cm}
    \begin{center}
		\includegraphics[width=1\textwidth]{../figures/raw/figure_F13_france_case_study.pdf}
	\end{center}
\end{figure}

